\chapter{Background Theory and Literature Review}
This chapter begins by examining the fundamental concepts in offline reinforcement learning and established methodologies for addressing distributional shift. Next, it explores diffusion models through the lens of variational bayesian framework, before unifying these concepts through the lens of continuous flows to introduce flow matching. The chapter concludes with an analysis of diffusion applications in trajectory synthesis and state-of-the-art approaches that integrate these models directly into reinforcement learning frameworks. Particular emphasis is placed on the mathematical formulations of diffusion and flow matching, as this constitutes the central focus of this dissertation and the subsequent methodologies.


\section{Offline Reinforcement Learning}
As mentioned in chapter 1, a significant hurdle in offline RL is distribution shift between the static dataset and the policy being learned. This results in standard off-policy RL methods to fail due to overestimation of Q-values, particularly when dealing with complex and multi-modal data distributions. When an agent's policy deviates from the offline dataset, it queries the Q-function for actions not present in the training data. This leads to erroneous and optimistically high value estimates. The problem is especially severe when the policy is chosen to maximize the Q-value, thereby exploiting regions of the state-action space where the Q-function is highly inaccurate. As a result, the Q-function believes it will perform exceptionally well, while actual performance is often poor.

\subsection{Policy Constraint Methods}
One early class of methods that’s been explored extensively is policy constraint methods. These approaches aim to mitigate distributional shift by updating the policy not just by maximizing the Q-value, but by doing so **su**bject to a constraint that the learned policy remains sufficiently close to the behavioral policy (collected data). This closeness is typically measured by a divergence metric, such as KL divergence or Maximum Mean Discrepancy (MMD).

To a degree, these methods can bound the distributional shift and the impact of out-of-distribution actions, potentially leading to more stable learning. Examples include **BEAR** (Bootstrapping Error Accumulation Reduction) by Kumar et al and **BRAC** (Behavior Regularized Actor Critic) by Wu et al [8], [9].

A significant challenge is their potential to be overly conservative, as a single constraint value $\epsilon$  might not work well across all states, especially if the behavior policy exhibits both highly random and highly deterministic actions. More critically, these methods often require **accurately estimating the behavior policy itself**, which can be extremely difficult, particularly when the data comes from diverse sources (e.g., humans or a mixture of different policies). Errors in this estimation can lead to an imperfect constraint and poor practical performance, making it a major bottleneck, especially during online fine-tuning.
\subsection{Implicit Policy Constraint Methods}
Building upon the insights into the limitations of explicit policy constraint methods, **implicit policy constraint methods**, such as **AWAC** (Advantage Weighted Actor Critic) by Nair et al. , emerged as an alternative that addresses the difficulty of behavior policy modeling. AWAC leverages a theoretical identity that shows the solution to a constrained optimization problem can be approximated using a **weighted maximum likelihood objective**. In this formulation, actions taken from the behavior policy are weighted by the exponentiated advantage function, effectively embedding the constraint implicitly [10]. More intuitively, AWAC approximates the best new policy by look at actions already in the dataset and give more weight (advantage) to the really good ones. 

This approach **obviates the need for explicit modeling of the behavior policy**, only requiring samples from it, which are readily available in the offline dataset. AWAC has shown to be effective for online fine-tuning after offline initialization and demonstrates strong performance on complex tasks like dexterous manipulation with human demonstration data.

While the sources highlight its advantages, it is still a form of policy constraint.  Its inherent conservativeness might limit the extent to which it can discover entirely novel, optimal behaviors far from the observed data.
\subsection{Model-based Offline RL Methods}
Taking a different algorithmic direction, **model-based offline RL methods**, exemplified by **MOPO** (Model-based Offline Policy Optimization) by Yu et al., propose to tackle the problem by learning a model of the environment's dynamics. While vanilla model-based RL (like MBPO) already showed some promise in offline settings, MOPO specifically modifies these algorithms to handle distributional shift. MOPO achieves this by **punishing the policy for exploiting out-of-distribution states** during model rollouts. This is implemented by subtracting an **uncertainty penalty** from the reward function, which becomes larger for state-action tuples that are more out of distribution (i.e., where the learned model's predictions are less certain) [11], [12].

MOPO hence has the ability to **generalize beyond the state and action support of the offline data**, enabling the learning of policies for tasks different from those for which the data was collected. It is more resilient to overestimation and overfitting than model-free methods and theoretically maximizes a lower bound of the policy's true return, balancing exploration and risk. Empirically, MOPO outperforms existing model-free offline RL methods and vanilla model-based approaches, especially on tasks requiring out-of-distribution generalization.

Nevertheless, models can still be inaccurate in regions far from the training data, and in a purely offline setting, these errors cannot be corrected through further interaction. The effectiveness relies heavily on the quality of the uncertainty estimation and the tuning of the penalty coefficient. It may also struggle with datasets lacking sufficient action diversity, and there are limits to its generalization capabilities if the new task requires exploring states completely outside the available data support.
\subsection{Conservative Q Learning}
Finally, a distinct and highly effective approach is **Conservative Q-Learning (CQL)** by Kumar et al., which directly addresses the core problem of **Q-value overestimation**. Instead of constraining the policy, CQL augments the standard Bellman error objective with a novel **Q-value regularizer**. This regularizer works by finding and minimizing (pushing down) erroneously high Q-values for out-of-distribution actions, while a more advanced version also compensates by pushing up Q-values for actions present in the dataset. This makes the learned Q-function inherently "pessimistic".

CQL is theoretically guaranteed to learn a **lower bound on the true Q-function**, effectively preventing the problematic overestimation. This direct approach leads to state-of-the-art performance across a wide range of tasks on standard benchmarks like D4RL and Atari. It significantly outperforms many previous offline RL methods, demonstrating the ability to "stitch" together parts of different trajectories to find more optimal paths, and achieving up to five times better results on challenging dexterous manipulation tasks. While conservative, the full CQL method only slightly underestimates Q-values, making it a robust and reliable algorithm.

The effectiveness of CQL depends on the appropriate selection of the regularization weight (alpha). Simpler variants of the algorithm without the compensatory term for in-distribution actions can lead to excessive underestimation. Despite its strong performance, ongoing research continues to explore further improvements.

In conclusion, policy constraint methods and implicit policy constraint methods both try to constrain the learned policy to remain close to the behavior policy. However, they do not aim to learn a pessimistic Q-function or ensure a lower bound on Q-values. The simply avoid the problematic evaluation of out-of-distribution action. On the other hand, MOPO and CQL are methods that embrace a conservative and “pessimistic” strategy in order to mitigate the risks associated with going out of distributions and value overestimation by either directly pushing down the Q-values or penalized rewards based on uncertainty. This ensures the learned policies are robust, accepting the trade-off that may lead to sub-optimal but reliable performance.

\section{Diffusion}
\section{Flow Matching}
\section{Planning with Diffusion}
\section{Other Relevant Works}


Note that references work like this \cite{wall} \cite{rtems}

\section{Summary}
Your final section should include a summary of the chapter, which will summarise what has been covered and where there are gaps in knowledge which pave the way into the original contributions you make in the coming chapters. 